\documentclass[10pt,a4paper]{amsart}
\usepackage[margin=19mm]{geometry}
\usepackage{amsmath,amssymb,mathtools}
\usepackage[scr=rsfs,cal=euler]{mathalfa}
\usepackage{xfrac}
\usepackage[unicode,hidelinks,bookmarks=false]{hyperref}
\newcommand{\ie}{i.e.\ }
\newcommand{\cf}{cf.\ }
\newcommand{\resp}{respectively\ }
\newcommand{\Subsection}{\subsection}
\providecommand{\nobreakdash}{\nobreak\hbox{-}}
\setlength{\emergencystretch}{2em}
\multlinegap0pt
\usepackage{mathtools}
\usepackage{amssymb}
\usepackage{enumerate}
\usepackage{tikz}
\newtheorem{lemma}{Lemma}
\newcommand\Z{\mathbb{Z}}
\newcommand\eps{\varepsilon}
\title{On the Real Zeroes of Half-integral Weight Hecke Cusp Forms}
\author{Jesse J\"a\"asaari}
\date{Source: arXiv:2409.15271v3 (CC BY 4.0)}
\begin{document}
\maketitle

\noindent\emph{Source and licence.} On the Real Zeroes of Half-integral Weight Hecke Cusp Forms. Version arXiv:2409.15271v3 (CC BY 4.0), distributed by arXiv under the Creative Commons Attribution 4.0 International licence, http://creativecommons.org/licenses/by/4.0/, as recorded in the arXiv licence metadata for this exact version and shown on the abstract page https://arxiv.org/abs/2409.15271v3. Full licence notice: \texttt{licenses/arxiv-2409.15271-CC-BY-4.0.txt}.\par\smallskip

\noindent\textit{Editorial scope.} Counted unit: the article's lemma final-key-lemma (source0.tex lines 1144--1148). The article's complete original proof of it is delivered with it (source0.tex lines 1150--1251, one \texttt{proof} environment of 102 lines). No mathematics is written, repaired or re-derived: the statement and the proof are the article's own lines, in the article's own notation, and no retained line is substituted or re-flowed.  Retained uncounted as the source-local context the proof needs: lines 205--207; lines 259--261; lines 540--544; lines 548--554; lines 582--586; lines 1104--1111.  Nothing was omitted from any retained span, so the package carries no omission record.  The retained lines cite 2 works, whose entries are transcribed verbatim from the article's own bibliography source in the e-print and delivered with short editorial labels recorded in the item's reference mapping.  No retained line contains a figure environment, an image-inclusion command or drawing code, so no image is delivered and the package declares no asset dependency.  Editorial formatting only, in addition: an AMS \texttt{amsart} preamble that restates the article's own declarations and macros used by the retained lines, namely the environments \texttt{lemma} on separate counters, together with the standard packages those lines need.  The retained environments print in this document as Lemma 1 in order of appearance, whereas the article prints the same items with the numbering of its own full document.  The complete native source of the article as distributed in the arXiv e-print for this exact version is bundled unchanged as \texttt{sources/165694/source0.tex}.
\begin{align}\label{Fourier-coeff-rel}
c_g(|d|p^2)=c_g(|d|)\left(\lambda_f(p)-\frac{\chi_d(p)}{\sqrt p}\right).
\end{align}

\begin{align}\label{connecting-normalisations}
\alpha_g|c_g(|d|)|^2=\omega_fL\left(\frac12,f\otimes\chi_d\right),
\end{align}

\begin{lemma}\label{large-sieve-2}
Let $v\geq 1$ be a fixed integer and $2\leq P<Q\leq 2P$. Then
\[ 
\sum_{f\in \mathcal B_k}\left|\sum_{P<p\leq Q}\frac{\lambda_f(p^v)}p\right|^2\ll_v k\frac1{P\log P}+k^{10/11}\frac{Q^{v/5}}{(\log P)^2}.\]
\end{lemma}

\begin{lemma}\label{Murty-Sinha}
Let $p$ be a prime. Then for any interval $[\alpha,\beta]\subset[-2,2]$ we have
\[ 
\frac{\#\{f\in \mathcal B_k:\,\lambda_f(p)\in[\alpha,\beta]\}}{\#\mathcal B_k}=\int\limits_\alpha^\beta\mathrm d \mu_p+O\left(\frac{\log p}{\log k}\right),\]
where
\[\mathrm d \mu_p:=\frac{p+1}\pi\cdot\frac{(1-x^2/4)^{1/2}}{\left(p^{1/2}+p^{-1/2}\right)^2-x^2}\,\mathrm d x.\]
\end{lemma}

\begin{lemma}\label{Halasz}
Let $f:\mathbb N\longrightarrow[-1,1]$ be a multiplicative function. Then we have
\[ 
\frac1X\sum_{n\leq X}f(n)\ll \exp\left(-\frac14\sum_{p\leq X}\frac{1-f(p)}p\right).\]
\end{lemma}

\begin{align}\label{non-zero-Fourier-coeff}
\sum_{k\sim K}\sum_{\substack{f\in \mathcal B_k\\
\omega_f L(1/2,f\otimes\chi_d)> k^{-1-2\theta}}}1 &\geq \sum_{k\in\Z}h\left(\frac{2k}K\right)\sum_{\substack{f\in \mathcal B_k \\
\omega_f L(1/2,f\otimes\chi_d)> k^{-1-2\theta}}}1 \nonumber\\
&>\left(\frac1{4}-\frac\eps2\right)\sum_{k\in\Z}h\left(\frac{2k}K\right)\sum_{f\in\mathcal B_k}1 \nonumber\\
&\geq\left(\frac14-\frac\eps2\right)(1-\nu)\#\mathcal S_K \nonumber \\
&\geq \left(\frac14-\eps\right)\#\mathcal S_K.
\end{align}

\begin{lemma}\label{final-key-lemma}
Let $\varepsilon>0$ be arbitrary but fixed, and $d\ll 1$ be a fixed fundamental discriminant\footnote{Recall that the function $h_g$ depends on $d$.}. Then there exists $X_0(\varepsilon)$ such that for $X_0<X<K$ we have, for all but at most $\varepsilon\# S_{K,d}$ forms $g\in S_{K,d}$,
\[ 
\left|\frac1X\sum_{n\sim X}|h_g(n)|-\frac1X\left|\sum_{n\sim X}h_g(n)\right|\right|\gg_{\eta,\varepsilon} 1.\]
\end{lemma} 

\begin{proof}

It suffices to show the following two assertions: 

\begin{enumerate}
\item For all but at most $(\varepsilon/2)\#S_{K,d}$ of the forms $g\in S_{K,d}$ we have 
\[ 
\frac1X\sum_{n\sim X}|h_g(n)|\gg_{\eta,\varepsilon}1,\]
where the implicit constant is independent of $g$. \\
\item For all but at most $(\varepsilon/2)\#S_{K,d}$ of the forms $g\in S_{K,d}$ we have
\[ 
\frac1X\sum_{n\sim X}h_g(n)=o(1).\]
\end{enumerate}
\noindent Towards (1) we argue as follows. Let us first assume that the fixed fundamental discriminant $d$ is positive. Then note that by the identity (\ref{Fourier-coeff-rel}), the triangle inequality, and Lemma \ref{Murty-Sinha} we have
\begin{align*}
\sum_{\substack{k\sim K\\
k\equiv 0\,(2)}}\sum_{\substack{g\in B_{k+\frac12}^+\\
\sqrt{\alpha_g}|c_g(|d|)|>k^{-1/2-\theta}}}\sum_{\substack{p\leq X\\
|c_g(|d|)^{-1}c_g(|d|p^2)|<p^{-\eta/2}}}\frac1p &\ll\sum_{\substack{k\sim K\\
k\equiv 0\,(2)}}\sum_{f\in \mathcal B_k}\sum_{\substack{p\leq X\\
|\lambda_f(p)-\chi_d(p)/\sqrt p|<p^{-\eta/2}}}\frac1p \\
&\leq \sum_{\substack{k\sim K\\
k\equiv 0\,(2)}}\sum_{f\in \mathcal B_k}\sum_{\substack{p\leq X\\
|\lambda_f(p)|<p^{-\eta/2}+1/\sqrt p}}\frac1p \\
&\ll\sum_{\substack{k\sim K\\
k\equiv 0\,(2)}} \#\mathcal B_k\sum_{p\leq X}\left(p^{-3/2}+p^{-1-\eta/2}+\frac{\log p}{p\log k}\right) \\
&\ll_\eta\sum_{\substack{k\sim K\\
k\equiv 0\,(2)}}\#\mathcal B_k\\
&\ll\sum_{\substack{k\sim K\\
k\equiv 0\,(2)}}\sum_{\substack{g\in B_{k+\frac12}^+\\
\sqrt{\alpha_g}|c_g(|d|)|>k^{-1/2-\theta}}}1,
\end{align*}
where in the last step we have used (\ref{non-zero-Fourier-coeff}) (together with (\ref{connecting-normalisations}) and the fact that $c_g(|d|)=0$ for positive $d$ when $k$ is odd), and the fact that $\#\mathcal B_k\asymp k$.

Hence there is a positive constant $C_1$ depending only at most on $\eta$ and $\theta$ such that for any given $\varepsilon>0$,
\begin{align*}
\sum_{\substack{p\leq X \\
h_g(p)=0}}\frac1p\leq\frac {C_1}{\varepsilon}
\end{align*}
for all but at most $(\varepsilon/2)\# I_{K,d}$ forms $g\in I_{K,d}$, where
\begin{align*}
I_{K,d}:=\bigcup_{\substack{k\sim K\\
k\equiv 0\,(2)}}\{g\in B_{k+\frac12}^+:\,\sqrt{\alpha_g}|c_g(|d|)|>k^{-1/2-\theta}\}.
\end{align*}
Similarly, for a fixed negative fundamental discriminant $d$, one shows that there is a positive constant $C_2$ depending only at most on $\eta$ and $\theta$ such that for any given $\varepsilon>0$,
\begin{align*}
\sum_{\substack{p\leq X \\
h_g(p)=0}}\frac1p\leq\frac{C_2}{\varepsilon}
\end{align*}
for all but at most $(\varepsilon/2)\# J_{K,d}$ forms $g\in J_{K,d}$, where
\begin{align*}
J_{K,d}:=\bigcup_{\substack{k\sim K\\
k\equiv 1\,(2)}}\{g\in B_{k+\frac12}^+:\,\sqrt{\alpha_g}|c_g(|d|)|>k^{-1/2-\theta}\}.
\end{align*} 
Note that for any given $d$, $S_{K,d}=I_{K,d}\cup J_{K,d}$ with one of the sets $I_{K,d},\, J_{K,d}$ being empty (as $c_g(|d|)=0$ when $(-1)^kd<0)$. Now the first assertion follows from \cite[Theorem 2]{Hildebrand1987}. 

For the second assertion, observe that, for a fixed positive fundamental discriminant $d$, 
\begin{align}\label{identity-towards-Halasz}
\sum_{\substack{p\leq X\\
h_g(p)=-1}}\frac1p=\sum_{\substack{p\leq X\\
c_g(|d|)^{-1}c_g(|d|p^2)<0}}\frac1p-\sum_{\substack{p\leq X\\
-p^{-\eta/2}<c_g(|d|)^{-1}c_g(|d|p^2)<0}}\frac 1p.
\end{align}
Recalling (\ref{Fourier-coeff-rel}) and using the fact that $\chi_d(p)\in\{-1,0,1\}$ it follows that the first sum on the right-hand side of (\ref{identity-towards-Halasz}) is
\begin{align*}
&\geq \sum_{\substack{p\leq X\\
\lambda_f(p)<-1/\sqrt p}}\frac1p \\
&\geq \sum_{\substack{p\leq X\\
\lambda_f(p)<0}}\frac1p-\sum_{\substack{p\leq X\\
-1/\sqrt p<\lambda_f(p)<0}}\frac1p.
\end{align*}
Arguing identically as in \cite[p. 1610]{Lester-Matomaki-Radziwill2018} we have that
\begin{align*}
\sum_{\substack{p\leq X\\
\lambda_f(p)<0}}\frac1p\geq\frac{1+o(1)}8\log\log X+\sum_{\log X\leq p\leq X^{1/1000}}\frac{\lambda_f(p^2)-2\lambda_f(p)}p
\end{align*} 
and using the large sieve inequality of Lemma \ref{large-sieve-2} the latter sum on the right-hand side contributes $o(\log\log X)$ for almost all forms $f\in\cup_{k\sim K\,\text{even}}\mathcal B_k$. 

To summarise, we have deduced that 
\begin{align*}
\sum_{\substack{p\leq X\\
h_g(p)=-1}}\frac1p\geq\frac{(1+o(1))}{8}\log\log X-\sum_{\substack{p\leq X\\
-1/\sqrt p<\lambda_f(p)<0}}\frac1p-\sum_{\substack{p\leq X\\
-p^{-\eta/2}<c_g(|d|)^{-1}c_g(|d|p^2)<0}}\frac 1p
\end{align*}
for almost all forms $g\in\cup_{k\sim K\,\text{even}}B_{k+1/2}^+$. But from the arguments used to establish the first assertion above we know that there is an absolute constant $C_3$ so that
\[ 
\sum_{\substack{p\leq X\\
-p^{-\eta/2}<c_g(|d|)^{-1}c_g(|d|p^2)<0}}\frac 1p\leq\frac{C_3}{\varepsilon}\]
for all but $(\varepsilon/4)\# I_{K,d}$ of the forms $g\in I_{K,d}$. A similar computation also shows that there exists an absolute constant $C_4$ so that
\[ 
\sum_{\substack{p\leq X\\
-1/\sqrt p<\lambda_f(p)<0}}\frac1p\leq\frac{C_4}{\varepsilon}
\]
for all but $(\varepsilon/4)\# I_{K,d}$ of the forms $g\in I_{K,d}$ (recall here that $f\in\mathcal B_k$ corresponds to some $g\in B_{k+\frac12}^+$ under the Shimura correspondence). Thus for all but at most $(\eps/2)\#I_{K,d}$ of the forms $g\in I_{K,d}$ we have
\[
\sum_{\substack{p\leq X\\
h_g(p)=-1}}\frac1p\geq\frac{(1+o(1))}{8}\log\log X. 
\]
Similarly one shows that the same estimate holds for all but $(\varepsilon/2)\# J_{K,d}$ of the forms $g\in J_{K,d}$ when the fixed fundamental discriminant $d$ is negative. This finishes the proof of (2) by Hal\'asz's theorem (Lemma \ref{Halasz}) by recalling that one of the sets $I_{K,d},\,J_{K,d}$ is empty for any given $d$. This concludes the proof of the lemma. 

\end{proof}

\begin{thebibliography}{99}
\bibitem[Hildeb]{Hildebrand1987} Adolf Hildebrand. \newblock Quantitative mean value theorems for nonnegative multiplicative functions. {II}. \newblock {\em Acta Arith.}, 48(3):209--260, 1987.
\bibitem[Lester]{Lester-Matomaki-Radziwill2018} Stephen Lester, Kaisa Matom\"{a}ki, and Maksym Radziwi\l\l. \newblock Small scale distribution of zeros and mass of modular forms. \newblock {\em J. Eur. Math. Soc. (JEMS)}, 20(7):1595--1627, 2018.
\end{thebibliography}
\end{document}
